% Part: turing-machines 
% Chapter: undecidability
% Section: introduction

\documentclass[../../../include/open-logic-section]{subfiles}

\begin{document}

% SEGMENT OLP-0265-S01

\olfileid{tur}{und}{int} 
\olsection{அறிமுகம்}

ஒவ்வொரு சார்பும், ஒவ்வொரு எண்கணிதச் சார்பும்கூடக் கணிக்கத்தக்கதாக
இருக்க முடியாது என்பது வெளிப்படையாகத் தோன்றலாம். எண்ணிலடங்காத பல
சார்புகளின் நடத்தை மிகவும் சிக்கலானது. எடுத்துக்காட்டாக, கதிரியக்கத்
துகள்களின் சிதைவிலிருந்து அல்லது வேறு குழப்பமான அல்லது எழுமாறான
நடத்தையிலிருந்து வரையறுக்கப்படும் சார்புகள் இத்தகையவை. ஒரு
குறிப்பிட்ட நேரத்திலிருந்து ஒரு நொடி இடைவெளிகளை எண்ணத் தொடங்கி,
அந்தத் தொடக்கக் கணத்துக்குப் பின்னான $n$-ஆம் ஒரு நொடி இடைவெளியில்
அண்டத்தில் சிதைவுறும் துகள்களின் எண்ணிக்கையாக $f(n)$ என்ற சார்பை
வரையறுக்கிறோம் என்று கொள்க. நாம் ஒருபோதும் கணிக்கலாம் என்று நம்ப
முடியாத ஒரு சார்புக்கான வேட்பாளராக இது தோன்றுகிறது.

எனினும், அத்தகைய சார்புகளை ஒருவர் எவ்வாறு கணிப்பார் என்று
கற்பனைசெய்ய முடியாதிருப்பது ஒன்று; அவை கணிக்கவியலாதவை என்று உண்மையில்
நிறுவுவது முற்றிலும் வேறொன்று. நம்பிக்கையளிக்காத அளவுக்குச் சிக்கலாகத்
தோன்றும் சார்புகள்கூட உண்மையில் ஒரு கருத்தியல் பொருளில் கணிக்கத்தக்கவையாக
இருக்கலாம். எடுத்துக்காட்டாக, காலத்தில் அண்டம் முடிவுறுவது என்று கொள்க
---சில அண்டவியல் கோட்பாடுகள் முன்னுரைப்பதைப்போல், தொலை எதிர்காலத்தில்
ஒருநாள் அண்டம் ஒரே புள்ளியாகச் சுருங்கும். அப்போது அந்தத் தொடக்கக்
கணத்திலிருந்து $f(n)$ வரையறுக்கப்பட்டிருக்கும் நொடிகளின் எண்ணிக்கை
முடிவுறுவதாக மட்டுமே இருக்கும் (எனினும் அது நம்பமுடியாத அளவுக்குப்
பெரியது). முடிவுறு எண்ணிக்கையிலான உள்ளீடுகளுக்கு மட்டுமே
வரையறுக்கப்பட்ட எந்தச் சார்பும் கணிக்கத்தக்கது: வெளியீடுகளை ஒரே பெரிய
அட்டவணையில் பட்டியலிடலாம் அல்லது மிகப் பெரிய ஒரு டியூரிங் பொறி
நிலைமாற்று வரைபடத்தில் குறியீடாக்கலாம்.

ஆம்/இல்லை கேள்விகளுக்கான விடைகளை மதிப்புகளாகத் தரும் சார்புகளின்
சிறப்பு நேர்வுகளில் நாம் அடிக்கடி ஆர்வம் கொள்கிறோம். எடுத்துக்காட்டாக,
``$n$ ஒரு பகா எண்ணா?'' என்ற கேள்வி பின்வரும் சார்புடன் தொடர்புடையது:
\[
\fn{isprime}(n) = \begin{cases}
  1 & \text{$n$ பகா எண் எனில்}\\
  0 & \text{இல்லையெனில்.}
  \end{cases}
\]
தொடர்புடைய $1/0$-மதிப்புச் சார்பு விளைவுறக் கணிக்கத்தக்கது எனில்,
ஆம்/இல்லை கேள்வியை \emph{விளைவுறத் தீர்மானிக்கலாம்} என்று கூறுகிறோம்.

விளைவுறக் கணிக்க முடியாத சார்புகளோ, விளைவுறத் தீர்மானிக்க முடியாத
சிக்கல்களோ உள்ளன என்று கணிதமுறையில் நிறுவ, கணக்கீட்டுக்கான ஒரு
குறிப்பிட்ட மாதிரியைத் தேர்ந்தெடுத்து, அதனால் கணிக்க முடியாத சார்புகள்
அல்லது தீர்மானிக்க முடியாத சிக்கல்கள் உள்ளன என்று காட்டுவது அவசியம்.
எடுத்துக்காட்டாக, ஒவ்வொரு சார்பையும் டியூரிங் பொறிகளால் கணிக்க முடியாது
என்றும், ஒவ்வொரு சிக்கலையும் டியூரிங் பொறிகளால் தீர்மானிக்க முடியாது
என்றும் காட்டலாம். எல்லாச் சார்புகளையும் கணிக்க டியூரிங் பொறிகள்
போதிய திறனுடையவை அல்ல என்பதோடு, எந்த விளைவுறு செயல்முறையாலும் அவற்றைக்
கணிக்க முடியாது என்று முடிவுசெய்ய சர்ச்--டியூரிங் முன்வைப்பை
அப்போது பயன்படுத்தலாம்.

இத்தகைய எதிர்மறை முடிவுகளை நிறுவுவதற்கான திறவுகோல், டியூரிங்
பொறிகளுக்கே எண்களை ஒதுக்க முடியும் என்ற உண்மை. இதைச் செய்வதற்கான
மிக எளிய வழி அவற்றைப் பட்டியலிடுவது; டியூரிங் பொறிகளையும் அவற்றின்
நிரல்களையும் எழுதுவதற்கு ஒரு குறிப்பிட்ட வழியைத் தேர்ந்து, அவற்றை
முறைப்படி பட்டியலிடலாம். இதைச் செய்ய முடியும் என்று கண்டதும்,
எளிய எண்ணளவுக் கருத்துகளிலிருந்து டியூரிங்-கணிக்கவியலாத சார்புகள்
உள்ளன என்பது பின்வருகிறது: $\Nat$ இலிருந்து~$\Nat$ குச் செல்லும்
சார்புகளின் கணம் (உண்மையில், $\Nat$ இலிருந்து $\{0, 1\}$ குச்
செல்லும் சார்புகளின் கணமேகூட) !!{nonenumerable}; ஆனால் எல்லா டியூரிங்
பொறிகளையும் பட்டியலிட முடிவதால், டியூரிங்-கணிக்கத்தக்க சார்புகளின்
கணம் !!{denumerable} மட்டுமே.

முறையே கணிக்கவியலாதவை என்றும் தீர்மானிக்கவியலாதவை என்றும் நிறுவக்கூடிய
\emph{குறிப்பிட்ட} சார்புகளையும் சிக்கல்களையும் நாம் வரையறுக்கலாம்.
அத்தகைய ஒரு சிக்கல், \emph{நிற்றல் சிக்கல்} எனப்படுவது. டியூரிங்
பொறிகளின் கட்டளைகளைப் பட்டியலிடுவதன் மூலம் அவற்றை முடிவுறுமாறு
விவரிக்கலாம். ஒரு டியூரிங் பொறியின் அத்தகைய விவரிப்பை, அதாவது ஒரு
டியூரிங் பொறி நிரலை, மற்றொரு டியூரிங் பொறிக்கான உள்ளீடாக இயல்பாகவே
பயன்படுத்தலாம். எனவே மற்ற டியூரிங் பொறிகளைப் பற்றிய கேள்விகளைத்
தீர்மானிக்கும் டியூரிங் பொறிகளைக் கருதலாம். குறிப்பாக ஆர்வமூட்டும்
ஒரு கேள்வி: ``கொடுக்கப்பட்ட டியூரிங் பொறி உள்ளீடு~$n$ இல்
தொடங்கப்படும்போது இறுதியில் நிற்கிறதா?'' இக்கேள்வியைத் தீர்மானிக்கக்கூடிய
ஒரு டியூரிங் பொறி இருந்தால் நன்றாக இருக்கும்: டியூரிங் பொறிகள்
முடிவிலிக் கண்ணிகளிலும் அவற்றைப் போன்றவற்றிலும் சிக்காமல் இருப்பதை
உறுதிசெய்யும் தரக்கட்டுப்பாட்டு டியூரிங் பொறியாக அதை எண்ணிக்கொள்க.
அத்தகைய டியூரிங் பொறி இருக்க முடியாது என்பதே டியூரிங் நிறுவிய
ஆர்வமூட்டும் உண்மை. ஒரு டியூரிங் பொறி~$M$ இன் விவரிப்பையும் ஏதோ ஓர்
எண்~$n$ ஐயும் கொண்ட உள்ளீட்டில் தொடங்கப்படும்போது,~$M$ ஆனது உள்ளீடு
$n$ இல் தொடங்கப்பட்டிருந்தால் நின்றிருக்குமா இல்லையா என்பதைப் பொறுத்து
வெளியீடு $1$ அல்லது $0$ உடன் எப்போதும் நிற்கும் ஒரே டியூரிங் பொறி
இருக்க முடியாது.

குறிப்பிட்ட தீர்மானிக்கவியலாத சிக்கல்களுக்கான எடுத்துக்காட்டுகள் நம்மிடம்
கிடைத்ததும், மற்ற சிக்கல்களும் தீர்மானிக்கவியலாதவை என்று காட்ட
அவற்றைப் பயன்படுத்தலாம். எடுத்துக்காட்டாக, ``முதல்தர
!!{formula}~$!A$ ஏற்புடையதா?'' என்ற கேள்வி புகழ்பெற்ற ஒரு
தீர்மானிக்கவியலாத சிக்கல். முதல்தர !!{formula}~$!A$ ஒன்றை
உள்ளீடாகக் கொடுக்கும்போது,~$!A$ ஏற்புடையதா இல்லையா என்பதைப் பொறுத்து
வெளியீடு $1$ அல்லது $0$ உடன் நிற்பது உறுதியான எந்த டியூரிங் பொறியும்
இல்லை. வரலாற்றில், இச்சிக்கலை விளைவுறத் தீர்ப்பதற்கான ஒரு
செயல்முறையைக் கண்டறியும் கேள்வி, வெறுமனே ``அந்தத்'' தீர்மானச்
சிக்கல் எனப்பட்டது; எனவே தீர்மானச் சிக்கல் தீர்க்கவியலாதது என்று
கூறுகிறோம். டியூரிங்கும் சர்ச்சும் ஏறத்தாழ ஒரே நேரத்தில் இம்முடிவைத்
தனித்தனியாக நிறுவினர்; எனவே இது சர்ச்--டியூரிங் தேற்றம் என்றும்
அழைக்கப்படுகிறது.

\end{document}
